% =============================================================================
% Results — Forecasting Models (modelling part)
% Prose chapter; the numeric tables are pinned --- do not edit the numbers.
% Numbers pinned from Report/metrics/model_comparison.md (generated by
% scripts/report/extract_model_comparison.py) and Report/metrics/run_timings.md.
% Test window 2025-10-24..2026-05-31, hourly (24p), paper-exact features.
% =============================================================================

\chapter{Results: Forecasting Models}
\label{ch:results-models}

This chapter reports the head-to-head comparison of the three forecasting
models specified in Chapter~\ref{ch:methods-models}. Six backtests were run,
covering LEAR, the DNN, and the Transformer, each on DK1 and DE, over a single
shared hourly test window (2025-10-24 to 2026-05-31) with the paper-exact
feature sets. Within each market the models are compared on the test days that
all of them produced. Because the Transformer's fourteen-day input sequence
consumes more leading days than the other pipelines, this common set is 187 days
on DK1 and 176 on DE. The chapter reports headline accuracy
(\S\ref{sec:rm-headline}), statistical significance (\S\ref{sec:rm-significance}),
and computation cost (\S\ref{sec:rm-cost}), then sets the results against the
published benchmarks (\S\ref{sec:rm-papers}) and closes with the modelling-side
threats to validity (\S\ref{sec:rm-validity}) and future work
(\S\ref{sec:rm-future}).

% -----------------------------------------------------------------------------
\section{Headline accuracy}
\label{sec:rm-headline}

\begin{table}[h]
  \centering
  \caption{DK1: final benchmark metrics (test 2025-10-24--2026-05-31, common days).}
  \label{tab:results-dk1-headline}
\begin{tabular}{lrrrr}
\toprule
 & MAE & RMSE & sMAPE & rMAE \\
\midrule
LEAR & 17.574 & 26.891 & 29.048 & 0.469 \\
DNN & 17.960 & 25.634 & 29.339 & 0.482 \\
Transformer & 20.076 & 28.909 & 31.767 & 0.543 \\
Naive (weekly) & 36.899 & 51.385 & 50.110 & 1.000 \\
\bottomrule
\end{tabular}
\end{table}

\begin{table}[h]
  \centering
  \caption{DE: final benchmark metrics (test 2025-10-24--2026-05-31, common days).}
  \label{tab:results-de-headline}
\begin{tabular}{lrrrr}
\toprule
 & MAE & RMSE & sMAPE & rMAE \\
\midrule
LEAR & 13.069 & 25.041 & 22.107 & 0.376 \\
DNN & 13.476 & 22.544 & 22.763 & 0.388 \\
Transformer & 14.719 & 23.755 & 23.901 & 0.420 \\
Naive (weekly) & 35.049 & 51.914 & 46.090 & 1.000 \\
\bottomrule
\end{tabular}
\end{table}

Tables~\ref{tab:results-dk1-headline} and~\ref{tab:results-de-headline} report
the four accuracy metrics on the common test days. The ordering is the same in
both markets and across MAE, rMAE, and sMAPE. LEAR is the most accurate, the DNN
follows closely, and the Transformer is third, with all three improving
substantially on the weekly-naive baseline. The rMAE values are 0.469 (LEAR),
0.482 (DNN), and 0.543 (Transformer) on DK1, and 0.376, 0.388, and 0.420
respectively on DE. Because rMAE normalises each model against the same-market
naive error, these figures are comparable across markets. All three models sit
markedly further below 1 on DE than on DK1, indicating that the German price was
the more predictable of the two relative to its own naive benchmark over this
window.

The single metric that reorders the models is RMSE. The DNN attains the lowest
RMSE in both markets (25.6 against LEAR's 26.9 on DK1, and 22.5 against 25.0 on
DE), and on DE LEAR in fact posts the highest RMSE of the three, behind even the
Transformer. As RMSE weights large deviations more heavily than MAE, this is
evidence that the DNN suppresses the largest forecast errors more effectively
than LEAR even on the days where LEAR's mean error is slightly lower. LEAR's
advantage is thus specific to the absolute-error metrics that Lago et
al.~\cite{lago2021}, \S5.4, designate as primary. On the squared-error criterion
the DNN is preferable. Whether either difference is statistically meaningful is
the subject of \S\ref{sec:rm-significance}.

% -----------------------------------------------------------------------------
\section{Statistical significance}
\label{sec:rm-significance}

\begin{table}[h]
  \centering
  \caption{DK1: GW multivariate p-values (column model better than row model).}
  \label{tab:results-dk1-gw}
\begin{tabular}{lrrrr}
\toprule
 & LEAR & DNN & Transformer & Naive (weekly) \\
\midrule
LEAR & -- & 0.836 & 1.000 & 1.000 \\
DNN & 0.164 & -- & 1.000 & 1.000 \\
Transformer & 0.001 & 0.001 & -- & 1.000 \\
Naive (weekly) & 0.000 & 0.000 & 0.000 & -- \\
\bottomrule
\end{tabular}
\end{table}

\begin{table}[h]
  \centering
  \caption{DE: GW multivariate p-values (column model better than row model).}
  \label{tab:results-de-gw}
\begin{tabular}{lrrrr}
\toprule
 & LEAR & DNN & Transformer & Naive (weekly) \\
\midrule
LEAR & -- & 0.807 & 0.988 & 1.000 \\
DNN & 0.193 & -- & 0.994 & 1.000 \\
Transformer & 0.012 & 0.006 & -- & 1.000 \\
Naive (weekly) & 0.000 & 0.000 & 0.000 & -- \\
\bottomrule
\end{tabular}
\end{table}

The Giacomini--White multivariate test~\cite{giacomini2006gw, lago2021}
establishes whether these differences are significant.
Tables~\ref{tab:results-dk1-gw} and~\ref{tab:results-de-gw} give the pairwise
p-values, where cell $(r,c)$ is the probability of a loss differential at least
as large as the one observed under the null that column model $c$ is no more
accurate than row model $r$ (the convention of~\cite{lago2021}, Fig.~5). A value
below 0.05 means that $c$ significantly outperforms $r$.

The central result is that LEAR and the DNN are statistically indistinguishable
on both markets. The hypothesis that LEAR is no more accurate than the DNN
cannot be rejected (DK1 $p=0.16$, DE $p=0.19$), and neither can its converse.
LEAR's numerical lead in the headline tables therefore does not amount to a
significant accuracy advantage, and the two models should be treated as
equivalent in accuracy. Both, however, significantly outperform the
Transformer (on DK1 at $p=0.001$ for each, and on DE at $p=0.012$ for LEAR and
$p=0.006$ for the DNN), and all three beat the weekly-naive baseline with $p\approx0$.
The comparison that survives significance testing is thus a two-way tie between
LEAR and the DNN at the top, a significantly weaker Transformer, and a baseline
that every model clears by a wide and significant margin, consistently across
DK1 and DE.

% -----------------------------------------------------------------------------
\section{Computation cost}
\label{sec:rm-cost}

\begin{table}[h]
  \centering
  \caption{Amortised recalibration cost (full-backtest wall-clock / common test day). Full timings in \texttt{Report/metrics/run\_timings.md}.}
  \label{tab:results-cost}
  \begin{tabularx}{\linewidth}{@{}l r r X@{}}
    \toprule
    \textbf{Model} & \textbf{DK1 s/day} & \textbf{DE s/day} & \textbf{Hardware / cadence} \\
    \midrule
    LEAR        & 50.3 & 41.3 & CPU, one LassoCV per period $\times$ 3 windows, daily \\
    DNN         & 0.9  & 1.6  & GPU, one MIMO net (single seed), daily \\
    Transformer & 1.6  & 1.9  & GPU, one MIMO net, weekly refit \\
    \bottomrule
  \end{tabularx}
\end{table}

Table~\ref{tab:results-cost} reports the amortised recalibration cost, the
full-backtest wall-clock divided by the number of common test days, for each
model and market. The figures invert the cost ordering of Lago et al.~\cite{lago2021},
Table~4, where on hourly CPU hardware LEAR trains in seconds and the DNN in
minutes. Here the relationship is reversed. The GPU-trained MIMO networks
recalibrate in one to two seconds per day, whereas LEAR (which fits an
independent cross-validated Lasso for each of the 24 delivery periods across
three calibration windows, entirely on CPU) takes forty to fifty. The DNN is
therefore roughly 25 to 55 times cheaper per recalibration than LEAR while, as
\S\ref{sec:rm-significance} showed, being its statistical equal in accuracy.
This inversion is a property of the hardware rather than of the models or the
markets, since a GPU rewards the dense neural networks and penalises LEAR's many
small per-period fits. On this platform it is the decisive practical
difference between two otherwise-tied models. The DNN buys equal accuracy at
around a fortieth of the compute. The Transformer, refitted weekly, carries a
comparably low amortised cost but does not convert it into competitive accuracy.

% -----------------------------------------------------------------------------
\section{Comparison with published benchmarks}
\label{sec:rm-papers}

A direct comparison of error levels with the source papers is complicated by the
shift in price regime. The benchmarks of Lago et al.~\cite{lago2021} were struck
on 2016--2018 data, whereas this study covers 2025--2026, so only rMAE levels and
rankings are comparable, never absolute errors (\cite{lago2021}, \S5.4).
Table~\ref{tab:results-vs-paper} sets the ensemble rMAE values against the
paper's Table~3 figures for the corresponding markets.

\begin{table}[h]
  \centering
  \caption{Our ensemble rMAE vs published (\cite{lago2021} Table~3). Periods differ, rankings/levels only.}
  \label{tab:results-vs-paper}
  \begin{tabularx}{\linewidth}{@{}l r r r r@{}}
    \toprule
    & \textbf{DK1 ours} & \textbf{NP \cite{lago2021}} & \textbf{DE ours} & \textbf{EPEX-DE \cite{lago2021}} \\
    \midrule
    LEAR & 0.469 & 0.420 & 0.376 & 0.433 \\
    DNN  & 0.482 & 0.407 & 0.388 & 0.374 \\
    \bottomrule
  \end{tabularx}
\end{table}

The levels are broadly in line with the published benchmark. On DK1 both models
are somewhat less accurate than the paper's Nord Pool figures (0.469 against
0.420 for LEAR, 0.482 against 0.407 for the DNN), whereas on DE the LEAR ensemble
is actually more accurate than the paper's EPEX-DE result (0.376 against 0.433)
and the DNN is marginally above it (0.388 against 0.374). That the absolute rMAE
levels land in the same range as a benchmark struck on entirely different data is
a modest validation of the replication in its own right.

The rankings, however, invert both source papers. Lago et al.~\cite{lago2021}
report the DNN as at least the equal of LEAR, and Llorente and
Portela~\cite{llorente2025transformer} report their Transformer ahead of the
DNN. The ordering here is the reverse on both counts, LEAR $\ge$ DNN $\ge$
Transformer. As \S\ref{sec:rm-significance} showed, the LEAR--DNN inversion is
only numerical, since the two are statistically tied, which is itself consistent
with Lago et al.'s observation that the gap is often insignificant. The
Transformer's third place, by contrast, is genuine and significant. The likeliest
cause is that the deep models are under-optimised relative to LEAR in this study
rather than inherently weaker, a point developed in \S\ref{sec:rm-validity}. The
DNN was given 30 tuning trials against the paper's 1500, and the Transformer reused the
paper's Nord Pool hyperparameters without per-market retuning and was adapted
into the rolling-refit pipeline rather than run in its native
train-once-and-renormalise mode (Chapter~\ref{ch:methods-models},
Table~\ref{tab:mm-deviations}), while LEAR re-selects its penalty by
cross-validation at every recalibration and so sits near its achievable best.

% -----------------------------------------------------------------------------
\section{Threats to validity (modelling side)}
\label{sec:rm-validity}

Several constraints bound the strength of these conclusions, all of them
stemming from data availability or compute budget rather than from the model
specifications. The test window is 31 weeks, short of the one-to-two years the
source papers recommend, because the usable price history begins only in March
2022 and the validation and calibration windows consume the remainder. Each
deep-model configuration is evaluated from a single backtest run at a fixed
seed. The DNN's four-seed averaging is applied only on the production serving
path and does not enter the reported backtest, and the Transformer is likewise a
single seed, so both deep models carry the seed variance of one draw, a caveat
that bears most heavily on the Transformer's third-place finish. The deep models are, as noted, tuned far less than LEAR
(Chapter~\ref{ch:methods-models}, Table~\ref{tab:mm-deviations}), so the rankings
should be read as contingent on the tuning budget. Finally, the 2025--2026
regime contains negative-price and high-renewable-infeed days largely absent from
the papers' data, and both fifteen-minute markets are downsampled to an hourly
mean for comparability. These choices affect every model identically and so do
not bias the paired tests, but they do limit comparability with the original
studies.

% -----------------------------------------------------------------------------
\section{Future work (modelling side)}
\label{sec:rm-future}

The most direct extension is to give the deep models a tuning budget
commensurate with LEAR's effective one before any architectural conclusion is
drawn. That means a full-scale TPE search for the DNN, approaching the paper's
1500 trials, and a per-market hyperparameter search for the Transformer in place
of the reused Nord Pool defaults. Beyond tuning, three further directions follow
from the deviations documented in Chapter~\ref{ch:methods-models}. The first is
to run the Transformer in its native train-once-and-renormalise
mode~\cite{llorente2025transformer} rather than the calibration-window refit used
here. The second is to restore the 1456-day LEAR window and the LARS+AIC penalty
selection of the original specification once the data history is long enough to
support them. The third is to extend the point forecasts to probabilistic ones
evaluated by pinball loss or CRPS. A per-hour
Diebold--Mariano profile, in place of the single daily-aggregate statistic, would
further localise where in the day the models diverge.

% =============================================================================
% END DRAFT
% =============================================================================
